\section{Method}
\label{sec:method}

\paragraph{Notation.}
Print size $s$ is the longest side in centimetres.
Ink spread $\sigma$ is an isotropic Gaussian standard deviation in millimetres.
Needle width is $w_n{=}0.30$\,mm and the tight-gap width is $g{=}0.50$\,mm.
Pairwise accuracy is $A$.
The effect is $A-0.5$, and $h$ is Cohen's $h$ of $A$ against $0.5$.
Skeleton agreement is $F_1$ at a 0.2\,mm tolerance.
The same symbols are used in the text, the tables and the captions.

\subsection{Physical rescale}
Every probe reads an RGB image, converts it to grey, and rescales it so the longest side spans the chosen print size $s$ at 10\,px/mm.
Downscaling uses area interpolation.
Upscaling uses cubic interpolation.
Sizes 3, 5, 10 and 15\,cm are the grid.
Preference tables use the 5\,cm probes.
The degradation table also records 10\,cm probes in the raw file; the paper reports the 5\,cm slice.

Ink is the set of pixels darker than an Otsu threshold \citep{otsu1979}.
If the grey standard deviation is at most 1, the threshold is fixed at mid-grey.
Line width is twice the distance transform on the ink mask, read on the Zhang--Suen skeleton and converted to millimetres \citep{zhang1984thinning}.

\subsection{Pre-registered probes}
Signs were fixed before any preference accuracy was computed.
$+1$ means a larger value is treated as more tattooable.
\texttt{line\_w\_p10\_mm} ($+1$) is the 10th percentile of skeleton line width.
\texttt{frac\_sub\_needle} ($-1$) is the fraction of skeleton pixels thinner than $w_n$.
\texttt{open\_ends\_per\_cm} ($-1$) counts skeleton endpoints per centimetre of skeleton.
\texttt{frac\_tight\_gap} ($-1$) is the fraction of the background skeleton inside the ink box that is narrower than $g$.
\texttt{midtone\_frac} ($-1$) is the fraction of pixels strictly between grey 60 and grey 200.
\texttt{border\_std} ($-1$) is the standard deviation of a 5\,px border.
\texttt{edge\_density\_per\_mm2} ($-1$) counts Canny edge pixels per square millimetre, with hysteresis 50 and 150 \citep{canny1986}.
\texttt{specks\_per\_cm2} ($-1$) counts connected ink components smaller than a $w_n$ square, per square centimetre.
An unweighted composite, the mean of signed robust $z$-scores, was defined in advance and is deprecated.
Appendix~\ref{app:probes} keeps one sentence on it.
Per-dimension accuracies in Table~\ref{tab:agree} are the probe result we interpret.
Edge transition (below) is excluded from the pre-registered feature list and included only in the ``all probes'' fusion.

\subsection{Post-hoc edge transition}
\texttt{edge\_transition\_mm} (sign $-1$) estimates how wide the ink-to-paper ramp is.
On Canny pixels we take the median Sobel magnitude and divide the median paper-minus-ink contrast by that magnitude, then convert pixels to millimetres.
A soft or blurred edge scores higher.
This probe was added after the first degradation run showed that blur and detail loss were invisible to the pre-registered set.
Section~\ref{app:prereg} records that order.
We do not treat its preference accuracy as a confirmatory test.

\subsection{Scores we compare against}
CLIPScore, PickScore, HPSv2.1 and ImageReward are used as released scorers \citep{hessel2021clipscore,kirstain2023pickapic,wu2023hpsv2,xu2023imagereward}.
Higher is treated as the model's preferred image.
We do not retrain them.
HPSv2.1 was trained on HPDv2, and ImageReward was trained on ImageRewardDB.
Accuracy of those two models on their own training distribution is leakage, not an external test.
PickScore was trained on Pick-a-Pic v1, so its accuracy on our Pick-a-Pic v2 tattoo-subject pairs is also leaky.
HPSv2.1 and ImageReward were not trained on Pick-a-Pic.
On HPDv2, every reward score covers the same tattoo-subject pairs as the probes.

Pairwise accuracy $A$ gives the human winner one point when its score is strictly higher, and half a point on a tie.
Intervals are 95\% prompt-cluster bootstrap percentiles from 2000 resamples with seed 0.
Each resample draws prompts with replacement and keeps every pair inside a drawn prompt, so the unit is the prompt.
A two-sided bootstrap $p$ is then Holm-adjusted within each dataset across all scorers in Table~\ref{tab:agree}.
A cell that sits on the smallest resolvable $p$ is reported as $p_{\mathrm{Holm}}\le\HpdHpstwoPholm$.
The fusion model is L2 logistic regression on the difference of standardised features, with GroupKFold by prompt at $k{=}5$ and seed 0.
Within each fold the winner side is randomly flipped so the label is balanced.
Held-out accuracy counts a pair as correct when the predicted probability that the original orientation is the human order exceeds one half.
Coefficients quoted in the text are the signs of a final model fit on all pairs.
Their magnitudes are in the released coefficient file, not restated as new claims.

\subsection{Controlled defects}
On \nDegrDesigns\ designs we apply six operators at two strengths, always from the clean grey image.
Thinning uses a dilation of the grey image with a $3$ or $5$ pixel square kernel (dark ink shrinks).
Thickening uses the matching erosion.
Haze blends the image toward a flat grey at weights 0.30 and 0.50 and adds a small Gaussian pixel noise.
Speckle sets a random fraction 0.003 or 0.01 of pixels to black.
Blur is a Gaussian with $\sigma$ of 1.5 or 3.0 pixels on the original raster, before physical rescale.
Detail loss downsamples by 4 or by 8 with area interpolation and nearest-neighbour upsample.
Detection is the fraction of designs on which the clean image outscores the degraded image under the same rule as pairwise accuracy.
Reward models were run on a subset of \nDegrRewDesigns\ designs.
The table uses that subset for every column so the cells match.

\subsection{Ink-spread stress test}
The spread model rescales to print size $s$ and applies an isotropic Gaussian of spread $\sigma$.
The default grid adds no pixel noise.
An earlier setting added grey-level noise with pixel standard deviation 4 and made skeleton $F_1$ fall with print size even at zero spread; that decline was an artifact, and those four noise-on values are reported only as a contrast.
The grid is sizes 3, 5, 10 and 15\,cm and spreads from 0 to 0.60\,mm in steps of 0.05\,mm.
Skeleton $F_1$ compares the clean ink skeleton with the spread ink skeleton.
The reference and the simulation use the same Otsu threshold.
A skeleton pixel matches if it lies within 0.2\,mm of the other skeleton.
Gap survival is the fraction of background pixels inside the clean ink box that remain background after the spread.
This is a stress test of the probes' scale assumptions.
It is not calibrated to photographs of healed tattoos.
